\section{Related work}\label{sec:related}
\subsection{Traffic forecasting backbones and the role of calibration}
Graph WaveNet combines temporal modelling with learned spatial dependence \citep{wu2019}. Its relevance here is the separation between a forecasting backbone and the uncertainty layer attached to its outputs. A graph predictor can exploit relationships between sensors, yet this does not determine how prediction errors should be used when their reports arrive late. The FAR-GW adapter adds causal validity and age channels and a noncrossing quantile head. These are documented study modifications rather than claims about the original implementation. Table~\ref{tab:related-models} distinguishes the forecasting references from the uncertainty mechanisms evaluated on their frozen outputs.

DCRNN represents spatial propagation through diffusion operations within a recurrent forecasting architecture \citep{li2018}. It supplies a contrasting temporal backbone and the provenance of the real benchmark packages. In this study, its multihorizon quantile adaptation uses the same causal channels and reporting schedules as the other predictors. The purpose of including it is to examine whether calibration effects persist under a different residual-generating model. Source-paper scores cannot serve as the comparison because dataset partitions, valid-target masks and fault processes differ. Only predictions generated under the common study protocol enter the paired backbone analysis.

STID provides a relatively simple forecasting design built around spatial and temporal identity representations \citep{shao2022}. STAEformer instead uses adaptive spatiotemporal embeddings with a Transformer \citep{liu2023}. These references demonstrate that useful traffic representations need not rely on the same graph-convolution construction. Their inclusion broadens the residual shapes encountered by the calibrator. It does not make the present study an exhaustive forecasting benchmark: only four fitted architectures are evaluated, and each has a documented adaptation. The distinction between model implementation, training protocol and calibration evaluation remains essential when interpreting transfer across them.

\subsection{Missing observations and uncertainty calibration}
Missing-input reconstruction is a related but distinct objective. GRIN uses graph message passing to reconstruct incomplete multivariate series \citep{cini2022}. An imputation method addresses what value the forecasting model can consume when a channel is absent; a feedback-aware calibrator addresses which forecast errors have actually become available. The two can interact, because input reconstruction changes the subsequent residual distribution, but they cannot be substituted for each other. Our causal last-observation and training-median fills are a deliberately explicit input policy. GRIN is discussed as relevant literature, not included as an unevaluated numerical comparator.

Conformal prediction separates an underlying predictor from a calibration procedure, as explained by \citet{angelopoulos2021}. A central distinction is between an empirical coverage estimate and a guarantee derived under a specified sampling assumption. Traffic observations are temporally dependent, and artificial outages further change which outcomes reach the service. Consequently, a calibration construction cannot acquire distribution-free validity merely because it uses residual quantiles. We retain the terminology of calibration while identifying the proposed rule as an empirical weighted signed-residual method. Its nominal interval level, observed coverage and registered decision tolerance are reported separately throughout the results.

Conformalized quantile regression combines quantile modelling with conformal calibration to accommodate heteroscedastic responses \citep{romano2019}. This is relevant to the use of a predictive centre and scale, because residual magnitudes need not be comparable across locations or traffic regimes without normalisation. Our quantile head supplies that scale but is not itself treated as a calibrated uncertainty guarantee. The final interval comes from received signed residuals and a fallback archive. The mathematical formulation therefore distinguishes predictor quantiles from calibration quantiles, making clear which quantities are fitted before test replay and which change after causal feedback release.

\citet{barber2023} analyse conformal prediction beyond exchangeability, including weighted quantiles and distribution drift. Their work provides a useful caution for interpreting relevance weights. Emphasising recent or contextually similar observations can address changing distributions, but the validity of a specific weighting scheme depends on its own assumptions and construction. The exponential target-age and context weights used here are not presented as an instantiation of that paper's complete theorem. They express a testable engineering hypothesis about useful residual evidence. The full ablations are needed precisely because plausible weighting components may help one network and harm another.

\subsection{Sequential, multistep and relational inference}
ACI updates a control variable from sequential coverage errors \citep{gibbs2021}. This establishes an important alternative to keeping a calibration archive completely fixed. For our comparison, the ACI adapter is horizon specific and updates once per spatial feedback batch. It cannot update before a packet arrives and receives no update from a permanently lost target. These changes define the actual comparator rather than the original source protocol. A favourable or unfavourable paired result therefore concerns the implemented adaptation under common causal replay, not a refutation or confirmation of the source theorem.

Sequential predictive conformal inference develops adaptive prediction intervals for time series by exploiting sequential residual information \citep{xu2023}. Its relevance is the broader principle that temporal dependence can carry information about future uncertainty. That principle supports examining a residual archive beyond an unordered static collection. However, relevance of historical residuals and availability of their outcomes remain separate questions. We discuss sequential predictive inference to position the problem, without reporting comparative scores that were never generated. Adding it as an executable comparator would require a separately specified feedback-arrival adaptation and a new registered evaluation.

Multistep adaptive conformal inference accounts for the fact that the outcome of a forecast is unavailable until its target time \citep{hallberg2024}. This ordinary horizon delay should not be confused with an additional communication backlog. A fifteen-minute forecast can have its outcome physically realised after fifteen minutes yet delivered much later. Our replay records both target and release timestamps and permits permanent loss. The ACI implementation uses the multistep reference as its starting point but changes spatial batching and reporting availability. Table~\ref{tab:related-calibration} records this distinction without implying that the source method ignores ordinary forecasting delay.

CoRel uses relationships between correlated time series to calibrate uncertainty on top of a pretrained predictor \citep{cini2025}. This is the nearest directly evaluated relational comparator and prevents an overbroad claim that neighbouring series have not previously been used for uncertainty estimation. Its study adaptation follows the official architecture while changing partitions, multihorizon targets and residual availability to match our protocol. CoRel is evaluated only with FAR-GW caches. The comparison consequently isolates calibration on a common predictor; it neither reproduces every source setting nor supports conclusions about CoRel on untested forecasting backbones.

\subsection{Recent work on incomplete and impaired feedback}
\citet{chen2026} study post-training adaptive conformal prediction for incomplete time series. The previously audited method uses imputation and missing-pattern-conditioned calibration, followed by adaptation when the test outcome is observed. Its substantial overlap with our motivation is acknowledged: incomplete histories can require uncertainty adjustments beyond ordinary clean-input calibration. Our different evaluation target concerns the arrival or permanent suppression of outcome feedback, represented separately from predictor-input availability. The distinction is limited to the reviewed protocol; it is not a claim that the entire incomplete-data literature lacks delayed reporting or that our method identifies unseen outcomes.

The September 2026 preprint by \citet{elhalabi2026} explicitly analyses ACI under delayed feedback. It relates forecasting delay to the persistence of residual information and derives coverage results for its delayed recursion. This directly narrows our positioning: studying delay is not new in itself. The traffic contribution is the joint replay of correlated input disruption, delayed backlog and permanent outcome loss, together with locked controls and component evidence. The preprint is labelled as such because publication status matters. Its theoretical results are not used to justify our weighted pooling rule or selectively missing outcome process.

The May 2026 preprint by \citet{wang2026} studies corrupted coverage feedback through binary flips, with robust filtering and active compensation. This addresses unreliable feedback rather than assuming every received indicator is correct. Our principal interventions suppress or delay reports rather than flip their labels, so the corruption models differ. Nevertheless, both settings show why the feedback channel deserves explicit treatment. The present experiments do not establish performance under adversarially corrupted values. The comparison table records these different problem definitions instead of assigning unsupported universal capabilities or ranking methods outside their evaluated tasks.

\subsection{Positioning and evidence boundary}
Table~\ref{tab:related-models} summarises representation choices and implementation scope; Table~\ref{tab:related-calibration} summarises uncertainty objectives, feedback assumptions and whether a method was numerically evaluated here. These tables are literature maps, not score tables. Numerical comparisons appear only in the results and share forecasts and masks. SUMO supplies controlled simulation for examining reporting and physical events separately \citep{lopez2018}. Combining these resources permits a focused benchmark contribution, while the recent literature rules out a broad first-of-its-kind claim for adaptive, relational or delayed-feedback inference. The remaining research question is conditional traffic interval quality under the precisely specified causal protocol.

The distinction between literature coverage and experimental coverage is important for an applied manuscript. A discussion can recognise a relevant method without suggesting that it has been reproduced. Conversely, an implemented comparator needs enough detail to identify changes to the original protocol. We therefore separate published algorithmic ideas, project adapters and newly proposed residual weighting. The evidence supports only comparisons for which frozen predictions, causal schedules and paired outcomes exist. This approach also makes missing baselines visible: they are opportunities for additional research, not hidden entries assigned favourable or unfavourable scores. The architecture table and calibration table expose these boundaries, while later numerical tables retain the exact direction of each effect and its uncertainty rather than collapsing every result into a superiority label.
